% Section III: System design and implementation (target 1.30 pages with Table I, Fig. 1, Fig. 2).
% Sources: docs/hardware.md, configs/real.yaml (real-1 final, frozen), docs/decisions.md
% (real-1, grip-5, rec-1, rec-6, hw-1). Data format and ACT execution details: Section IV.
\section{System Design and Implementation}
\label{sec:system}

\subsection{Hardware and Workspace}
\label{sec:hardware}
The follower is a six-joint Rainbow Robotics RB5-850 collaborative
manipulator. A host computer sends commands to its control box and receives
the robot state over two TCP channels through the vendor's Python API.
The operator uses the right arm of an RB-Y1 leader with seven Dynamixel joints.
Its handle carries a clutch button, a mode button, and an analog trigger.
The leader is read without torque while the clutch is pressed. When the
clutch is released, its joints are torqued to hold their pose. The end
effector is an in-house 3D-printed parallel gripper driven by one Dynamixel
MX-106 actuator on a separate serial bus.

Three RGB cameras record at $640\times480$ and 30~fps: a RealSense D435
fixed obliquely in front of the workspace (\emph{front}), a Logitech C270
fixed above the item area and looking straight down (\emph{top}), and a
Logitech Brio~101 mounted on the gripper and aimed at the fingers
(\emph{grip}). Exposure, gain, and white balance are fixed. Each exposure is a
multiple of 1/120~s, so the 60~Hz room lighting does not flicker in the images.
At start-up, each camera's mean brightness is compared with a reference value,
and a deviation above 15\% requires operator confirmation.
Depth is not used. A black box and a meal box with a large and a small hole
are fixed at marked positions on the robot's table, and the items are placed
within the view of the top camera (Fig.~\ref{fig:setup}). Table~\ref{tab:hardware}
summarizes the components and operating settings.
% TODO: item dimensions, finger clearance, hole sizes, container positions in the base frame.

\begin{figure}[t]
\centering
\fbox{\parbox[c][0.55\columnwidth][c]{0.95\columnwidth}{\centering
Photograph placeholder: RB5-850 with gripper and grip camera, RB-Y1 leader,
front and top cameras, black box, meal box (large and small holes), item area.}}
\caption{Experimental setup. The front and top cameras are fixed; the grip
camera moves with the gripper.}
\label{fig:setup}
\end{figure}

\begin{table}[t]
\caption{Hardware, Software, and Operating Settings}
\label{tab:hardware}
\centering
\footnotesize
\setlength{\tabcolsep}{3pt}
\begin{tabular}{@{}lp{0.68\columnwidth}@{}}
\toprule
Item & Specification\\
\midrule
Follower & RB5-850, 6 joints; control box software 8.4.4\\
Leader & RB-Y1 right arm, 7 joints (XM540/XM430), 2~Mbps\\
Gripper & In-house parallel, MX-106, current-based position mode\\
Cameras & D435 (front), C270 (top), Brio~101 (grip); RGB $640\times480$, 30~fps\\
Host & Intel i7-13700H, NVIDIA RTX A1000 6~GB laptop GPU, Ubuntu 24.04\\
Software & rbpodo 0.16.14, LeRobot 0.6.1, PyTorch 2.11\\
Command loop & 50~Hz Cartesian targets (\texttt{move\_servo\_l})\\
Speed limits & 250~mm/s in $r$ and $z$, 360~mm/s along $\phi$;
J1 75~deg/s, J2--J6 60~deg/s; 2.5~m/s$^2$\\
Workspace & $r$ 250--800~mm (depends on $z$), $z$ 19--465~mm,
$\phi$ $-160$ to $133^\circ$\\
Gripper force & 0.50~A current cap while moving, 0.84~A while grasping\\
\bottomrule
\end{tabular}
% TODO: TCP offset, tool mass, and collision-detection setting used in the evaluation
% (the TCP offset on the controller is still provisional; docs/hardware.md).
\end{table}

\subsection{Software Architecture}
\label{sec:architecture}
All software runs on the host as Python processes, arranged in a recording
path and an autonomous-execution path that share the camera, gripper, and
command-shaping components (Fig.~\ref{fig:architecture}).

In the recording path, a teleoperation process runs a 50~Hz cycle.
In each cycle, it reads the leader joints and the latest RB5 state,
computes the follower target, and sends it.
The gripper runs in its own 50~Hz thread, so a gripper fault does not stop
the arm. The cameras run in a separate process with a capture thread and a
writer thread per camera, which store every frame as a JPEG image with its
capture time. Image work therefore cannot delay the command cycle. Robot,
gripper, and camera samples are stamped with the same monotonic clock and
saved at their native rates. The 30~Hz dataset is built from them offline
(Section~\ref{sec:workflow}).

In the execution path, a policy process opens the same cameras with the same
fixed settings and runs ACT on the GPU, at most once per 30~Hz step. Each
predicted action chunk is stamped with its observation time. A servo
process keeps the 50~Hz cycle without waiting for the policy. It combines the
chunks by time and passes the resulting target through the same command
shaping as in teleoperation. The two processes are pinned to separate CPU
cores. The stated rates are implementation targets, not hard real-time
guarantees. A cycle whose command would be sent more than 30~ms after the
previous one keeps the previous target, and three such cycles in a row stop
the run.

\begin{figure}[t]
\centering
\fbox{\parbox[c][0.62\columnwidth][c]{0.95\columnwidth}{\centering
Diagram placeholder.\\[3pt]
Recording: leader FK $\rightarrow$ clutch-relative cylindrical mapping $\rightarrow$
command shaping $\rightarrow$ \texttt{move\_servo\_l} (50~Hz); gripper thread;
camera process (30~fps JPEG); logs with one clock $\rightarrow$ 30~Hz dataset.\\[3pt]
Execution: cameras + robot state $\rightarrow$ ACT (policy process, GPU)
$\rightarrow$ time-stamped chunks $\rightarrow$ servo process (50~Hz) $\rightarrow$
same command shaping $\rightarrow$ \texttt{move\_servo\_l}; gripper thread.}}
\caption{Recording and autonomous-execution paths. Both send absolute
Cartesian targets through the same command shaping; the recorded action is
the shaped target.}
\label{fig:architecture}
\end{figure}

\subsection{Dissimilar Leader--Follower Interface}
\label{sec:interface}
The leader's joint angles are converted into the position of its wrist
center $\mathbf{p}_L=(x_L,y_L,z_L)$ by forward kinematics on the vendor's URDF
model, with $x_L$ pointing forward and $z_L$ up.
The seven leader joints and the six follower joints are never matched one to
one. The follower target is expressed in cylindrical coordinates
$(\phi, r, z)$ about the RB5 base axis: azimuth, horizontal distance from the
axis, and height. When the clutch is pressed at time $t_0$, the interface
stores $\mathbf{p}_L(t_0)$ and the current target $(\phi_0, r_0, z_0)$. Each
axis of the leader displacement $\Delta\mathbf{p}_L=\mathbf{p}_L(t)-\mathbf{p}_L(t_0)$
then drives a single follower coordinate. In the normal mode,
\begin{equation}
\phi = \phi_0,\quad r = r_0 + s\,\Delta x_L,\quad z = z_0 + s\,\Delta z_L,
\label{eq:normal}
\end{equation}
and while the mode button is held,
\begin{equation}
\phi = \phi_0 + k\,\Delta y_L,\quad r = r_0,\quad z = z_0,
\label{eq:rotate}
\end{equation}
with $s=1$ and $k=0.3^\circ$/mm. The commanded position is
$\mathbf{p}_d=(r\cos\phi,\ r\sin\phi,\ z)$ in the base frame. Releasing the
clutch freezes the target, and the next press re-anchors both references, so
the operator can reposition the hand without moving the robot.

Early teleoperation trials used a Cartesian mapping,
$\mathbf{p}_d=\mathbf{p}_d(t_0)+s\,\Delta\mathbf{p}_L$. In these trials, the
natural arcs of the operator's hand leaked into unintended follower axes, and
sideways drift turned the base. With \eqref{eq:normal} and
\eqref{eq:rotate}, forward and backward leader motion moves the gripper along
the arm's radial direction, vertical motion changes only its height, and the
base turns only in the rotate mode. A dead band
of 0.3~mm and a first-order low-pass filter with a 30~ms time constant
suppress encoder-step jitter of $\mathbf{p}_L$.

The tool orientation is not teleoperated. Roll and pitch keep the ready-pose
values, with the gripper pointing straight down. The yaw changes with the
target azimuth $\phi$, so the gripper turns with the arm and the follower's
last joint stays still. The orientation is thus a fixed function of the target
position, and the position together with a gripper value fully specifies the
command. Items must therefore be graspable without wrist rotation.

The gripper opening is proportional to the trigger. At start-up the gripper
follows only after the trigger reaches its current opening, which prevents
a jump. In current-based position mode, the motor current is limited to
0.50~A while moving. When the fingers stall on an item, the goal is set
60~ticks ($5^\circ$) past the stall point and the cap is raised to 0.84~A,
so the item is held firmly.

\subsection{Operating Limits and Robot Interface}
\label{sec:limits}
In every 20~ms cycle, the desired position passes through a fixed sequence of
command-shaping stages. Equations \eqref{eq:normal} and \eqref{eq:rotate}
give this position in teleoperation; a policy action gives it in autonomous
execution. The first two stages use the same cylindrical coordinates as
the mapping.
\begin{enumerate}
\item \emph{Workspace}: the target is clipped to
$r_{\min}(z)\le r\le r_{\max}(z)$, $19\le z\le465$~mm, and
$-160^\circ\le\phi\le133^\circ$ (Table~\ref{tab:hardware}). The lower
height bound keeps the fingers just above the table.
\item \emph{Speed}: the component of each step along $\phi$ is limited to
360~mm/s, and the component in $r$ and $z$ to 250~mm/s. The target may lag the
hand by up to 20~mm and catches up after the hand stops.
\item \emph{Singularities and joint limits}: the joint angles of the step are
predicted by the RB5's analytic inverse kinematics. Wrist, shoulder, and
elbow singularities and joint limits are monitored. Near one of them, the
speed limit falls linearly to 25\%; closer still, the step is not sent.
\item \emph{Joint speed}: per-joint limits of 75~deg/s for J1 and 60~deg/s
for J2--J6.
\item \emph{Acceleration}: a 2.5~m/s$^2$ limit on the target, applied
last. In trials, the largest accelerations occurred when an earlier limit
switched on or off.
\end{enumerate}
The shaped position and the orientation rule form the TCP pose that is
streamed to the controller with the RB5 API call \texttt{move\_servo\_l}.
In one teleoperation trial, the arm followed the commands about 80~ms
behind, and the joint error stayed within $0.23^\circ$. The speed limits
were raised step by step from 100~mm/s, based on hand speeds measured in
logged real-robot trials. They were then frozen together with all other
teleoperation settings before demonstration collection.

Before any motion, the software checks that the controller is in real mode
and that its TCP offset matches the expected value. It moves the robot to a
fixed ready pose at low joint speed only after the operator types a
confirmation. During operation, stale robot state, repeated leader read
failures, or repeated late cycles stop the robot.
The limiter is not collision avoidance and offers no formal safety
guarantee. No collision-aware planner is used, and the controller's
internal motion generation is outside the scope of this paper.
% TODO: state the controller collision-detection setting and tool mass used in the evaluation.
